\chapter{Literature Review}
\label{ch:literature-review}

\section{Introduction}
\label{sec:lr-intro}

This chapter provides a review of the published literature on the four pillars of NirmalNode, (i) the scale and health impact of industrial and ambient air pollution with special attention to Bangladesh, (ii) Internet of Things (IoT)-based air quality monitoring systems, (iii) artificial intelligence (AI)/machine learning (ML)-based air pollution prediction, and (iv) air purification technologies, specifically HEPA filtration, activated-carbon adsorption, and cyclonic pre-separation. The chapter also discusses the emerging concepts of Green IoT, Green AI and incremental/online (Adaptive) machine learning which together form the methodological backbone of the proposed system. Each section ends with a short synthesis, and the chapter concludes with a consolidated comparison table and a clear statement of the research gap that NirmalNode intends to fill.

\section{Industrial and Ambient Air Pollution: Global and Bangladesh Context}
\label{sec:lr-pollution-context}

Research on air pollution over the past decade has converged on the finding that exposure to fine particulate matter (PM2.5) is associated with health risks, even at concentrations previously deemed safe. The World Health Organization revised their 2021 guideline to lower the recommended annual limit of PM2.5 from 10 to 5 $\mu$g/m$^3$, and subsequent modeling work has shown that the \emph{non-anthropogenic} (natural) background component of PM2.5 already accounts for a substantial share of this reduced allowance in many regions, meaning that anthropogenic sources -- including industrial emissions -- must be reduced even further to remain within the new guideline \citep{pai2022who}. In a parallel analysis of small and medium cities, we conclude that the stricter WHO guideline is likely to be extremely difficult to meet in such cities without dedicated local-scale intervention, as they generally lack both the monitoring infrastructure and the regulatory capacity available to large metropolitan centers \citep{polezer2023new}.

A national evidence review in Bangladesh, in particular, finds vehicular emissions, brick kilns and unplanned industrial growth as the leading contributors to the country’s air pollution burden, and records its significant economic and public-health costs \citep{khatun2023evidence}. City-level and district-level studies support this at finer resolution. Long-term PM2.5 exposure in Bangladesh is reconstructed using satellite and ground-based data, revealing not a spatially uniform pattern of exposure, but rather specific and persistent pollution hotspots. These hotspots are linked to long-term health risk \citep{ali2025longterm}. In Dhaka, the Tejgaon industrial area shows significant variations in PM2.5 and PM10 values and corresponding health risk indices with distance from specific industrial sources \citep{basak2025spatial}, while a satellite-based (Sentinel-5 TROPOMI) trend analysis of Chattogram – the second-largest industrial hub of the country – reveals this city to be relatively less-studied than Dhaka despite a large pollutant burden \citep{huraira2024spatial}. Similarly, landfill-adjacent industrial zones in Dhaka have been found to emit and disperse gaseous pollutants (CH$_4$, CO$_2$ and other non-methane organic compounds) that had not previously been quantified in their atmospheric impact \citep{aurnab2026monitoring}. Remote-sensing analysis of atmospheric CO$_2$ anomalies over Bangladesh associates elevated concentrations with anthropogenic activity concentrated in and around industrial and urban centers \citep{ghosh2026spatiotemporal}. Industrial effluent near Dhaka has also been shown to increase trace-element concentrations in the surrounding environment to levels of ecological and health concern \citep{lisa2025investigative}, and general population studies in pollution-prone urban areas of Bangladesh report measurable self-reported health deterioration associated with perceived air pollution exposure \citep{rahman2025perceived}.

\textbf{Synthesis.} This work establishes two facts that directly motivate the present capstone: (1) the air pollution problem in Bangladesh is dominated by identifiable, spatially concentrated industrial and urban hotspots, rather than being spatially uniform, (2) existing evidence on these hotspots is drawn almost entirely from passive environmental monitoring (satellite remote sensing, periodic ground sampling), rather than from continuous, real-time, hotspot-embedded sensing capable of driving an automated response.

\section{Occupational Health Impacts at Industrial Hotspots}
\label{sec:lr-occupational-health}

A second branch of the literature reports the direct occupational health impacts of pollutant exposures targeted by NirmalNode. Elevated PM2.5 exposure in fiber-cement roofing manufacturing workers is associated with measurable impairment of lung function \citep{nasri2023pm25}. Exposure to particulate matter at construction sites has various associated negative health effects \citep{sekhavati2023particulate}. In addition to the documented respiratory and cardiovascular effects, reviews of PM2.5 exposure across multiple industry types report associations with mental health conditions, including elevated risk of depression and anxiety, and more severe psychiatric outcomes in some studies \citep{alhadhrami2024pm25mental}. Collectively, these findings indicate that the health burden of industrial air pollution goes beyond respiratory illness, affecting broader physical and psychological well-being, highlighting the urgency to protect workers at the specific points where their exposure is highest -- the very hotspots NirmalNode aims to address.

\section{IoT-Based Air Quality Monitoring Systems}
\label{sec:lr-iot-monitoring}

A large body of engineering literature has demonstrated the feasibility of constructing functioning air-quality monitoring networks using low-cost microcontrollers and commodity sensors. \citet{alam2025iot} suggest an IoT-based environmental monitoring system for real-time monitoring in developing areas, with a focus on low hardware cost and remote access. \citet{zhang2020realtime} combine mobile and fixed Internet of Things (IoT) sensing nodes to realize localized air quality monitoring and prediction, and show that a hybrid fixed/mobile sensing network can improve spatial resolution over fixed-only networks. \citet{kataria2022aiiot} propose a hybrid AI-IoT model for air quality prediction in a smart-city environment with network assistance and incorporate communication-network parameters into the prediction pipeline. Following a broadly similar sense-transmit-predict pipeline, \citet{mani2021aipowered} propose an AI-powered IoT system for real-time air pollution monitoring and forecasting. At system level, \citet{ramadhani2025iot} propose a low-cost-sensor based IoT implementation for air quality monitoring and \citet{pavan2025realtime} present a similar low-cost-sensor based IoT system with buzzer and SMS based alerting and ThingSpeak based remote visualization. More generally, a review of low-cost sensor networks for air quality management in developing countries contends that these networks are not merely a cheap substitute for reference-grade monitoring but a practical necessity given the financial and infrastructure constraints of these contexts \citep{shabbir2025paradigm}. Similarly, work on applying AI to model indoor air quality dynamics in developing nations, which complements this work, demonstrates that low-cost, field deployed sensing captures pollution patterns driven by cooking methods, ventilation and living density that are highly context-specific and poorly captured by generic models trained elsewhere \citep{karmakar2024exploring}. Similarly, a review chapter on AI-based air pollution monitoring in developing countries identifies data integration, infrastructure constraints, and energy-sector limitations as recurrent barriers to AI deployment at the policy scale \citep{noman2025aiair}. A combination of satellite analytics with AI and IoT sensing has also been proposed to support regulatory enforcement against industrial air pollution at a broader policy and enforcement level in the South Asian context \citep{ahmed2025smart}.

\textbf{Synthesis.} These systems together provide a mature and well-validated concept for IoT-based monitoring and even AI-assisted prediction. But in every case we looked at, the system’s role is limited to reporting or predicting a pollutant level – none of the IoT/AI monitoring systems we looked at are coupled to an automated physical localized mitigation action such as a purifier .

\section{Air Purification Technologies}
\label{sec:lr-purification}

\subsection{HEPA Filtration}
The High-Efficiency Particulate Air (HEPA) filtration is the most rigorously tested technology for removing fine particulate matter from air streams with a defining performance standard that at least 99.97\% of particles 0.3 $\mu$m in diameter are captured. Field evaluations demonstrate such effectiveness in real-world conditions: portable HEPA air cleaners have been shown to greatly reduce particulate concentrations in homes near major highways \citep{cox2018effectiveness}, and a controlled real-world evaluation of portable air cleaners confirms significant reductions in home particulate matter concentrations under typical operating conditions \citep{lu2023realworld}. The use of HEPA equipped purification systems installed in intercity busses has been shown to reduce the concentrations of particulate matter and airborne bacteria \citep{lee2022effect} and standalone HEPA air cleaners have been shown to significantly improve indoor PM2.5 concentrations in residential and clinical settings \citep{chen2022efficacy}. Portable HEPA filtration has shown the ability to substantially reduce indoor PM2.5 exposure even under severe pollution loading, e.g., wildland-urban-interface wildfire smoke intrusion \citep{chen2025fine}. We compare HEPA-based air purifiers to ionization-based and other purification technologies, and find that mechanical HEPA filtration is more reliably effective, especially for larger particle-size fractions, than ion-generator-based approaches \citep{dubey2021assessing}. Similarly, a broader review of indoor air purifier technologies for mitigating microbial contamination identified HEPA-based mechanical filtration as one of the most consistently effective methods evaluated \citep{szczotko2022review}, and focused microbiological-efficiency studies of HEPA filtration in medical-facility settings have reported comparably excellent performance for airborne biological contaminants \citep{laine2025microbiological}.

\subsection{Activated Carbon Filtration}
HEPA filters are effective against particulate matter but not at removing gaseous pollutants such as VOCs and odor causing compounds. The complementary technology for this purpose is the use of activated carbon filtration which works on the principle of adsorption of gas molecules onto the surface of the activated carbon extensive internal pore surface area. This is a key stage in the purification train proposed for NirmalNode.

\subsection{Cyclone Separation}
Cyclonic (centrifugal) separation is a well established technique used as a pre-filtration device. It induces a tangential swirling flow, which throws heavier particles outwards against the wall of the separator under centrifugal force, thus removing coarse particulate matter from an air stream before it reaches finer filtration stages. As proposed for NirmalNode, placing a cyclone separator before a HEPA filter extends the operational lifetime of the (more expensive) HEPA stage by removing the coarse particulate load before it reaches the filter media.

\subsection{Adoption Barriers}
These technologies have been shown to be effective, but their adoption is very low in Bangladesh. A multi-phase field study of household air purifier adoption in Dhaka finds that fewer than 1\% of middle-class households that can afford a purifier actually own one. This is driven substantially by misbeliefs about relative indoor/outdoor air quality and about purifier effectiveness, resulting in a willingness to pay far below the retail cost of a unit \citep{chowdhury2025misbeliefs}.

\textbf{Synthesis} The literature on purification substantiates that HEPA and activated-carbon filtration, aided by cyclonic pre-separation, are well-established technologies for the removal of both particulate and gaseous pollutants. All of the purification studies reviewed evaluate these technologies as \emph{room-scale or vehicle-scale} appliances reacting on the current air conditions, rather than as a compact, worker-proximate, hotspot-targeted unit operating on a predictive (1--2 hour look-ahead) activation schedule (which is the specific configuration proposed in this capstone).

\section{Green IoT and Green AI}
\label{sec:lr-green}

Green IoT is the design of IoT hardware and communication systems that reduce energy consumption, material use, and computational overhead while maintaining functional performance, for instance by using low-power microcontrollers, duty-cycled sensing and efficient wireless protocols. Green AI applies the same thinking to the machine-learning layer: instead of defaulting to big, computationally expensive deep-learning models, Green AI promotes lightweight models that are accurate enough for the task at hand, while dramatically reducing training time, inference latency, memory footprint and energy draw. The philosophy is implicit rather than explicit in the air-quality literature reviewed: systems such as those of \citet{alam2025iot} and \citet{shabbir2025paradigm} are motivated primarily by cost constraints in developing-country deployment, which aligns closely with, but is not formally framed as, the Green IoT/Green AI paradigm that this capstone adopts explicitly as a design principle.

\section{Adaptive AI and Incremental (Online) Learning}
\label{sec:lr-adaptive-ai}

None of the AI-based air-quality prediction systems reviewed in Section~\ref{sec:lr-iot-monitoring} reports using incremental or online learning. Instead, each system trains a model once, on a fixed historical dataset, and then deploys it statically. This is a serious limitation for any system to be deployed en masse over structurally different sites, since a model trained on the pollution signature (traffic patterns, industrial mix, meteorology) of one location cannot be expected to generalize perfectly to another site without some form of retraining or adaptation. Incremental (online) learning algorithms, e.g., SGD Regressor, Passive-Aggressive Regressor, Hoeffding Trees, and Adaptive Random Forests (several of them are implemented in the open-source River library), update model parameters continuously as new labeled data arrives without requiring the full dataset to be reprocessed. Such algorithms are well established in the broader literature on online learning, but to the best of the author’s knowledge from reviewing the supplied reference set, have not previously been applied to industrial, hotspot-level air pollution prediction coupled to an automated purification response. Chapter 3 discusses NirmalNode’s Adaptive AI design, the central methodological contribution that fills this gap in the literature.
\section{Comparative Analysis of Reviewed Systems}
\label{sec:lr-comparison}

Table~\ref{tab:lr-comparison} summarizes the reviewed IoT/AI-based air-quality systems along the dimensions most relevant to NirmalNode: sensing scope, prediction capability, purification capability, and adaptivity of the underlying model.

\begin{table}[htbp]
\centering
\caption{Comparative summary of reviewed IoT/AI-based air quality systems}
\label{tab:lr-comparison}
\small
\begin{tabular}{@{}p{2.6cm}p{2.6cm}p{2.1cm}p{2.6cm}p{2.6cm}@{}}
\toprule
\textbf{System / Study} & \textbf{Sensing Scope} & \textbf{Predictive?} & \textbf{Purification Coupled?} & \textbf{Model Adaptivity} \\
\midrule
\citet{alam2025iot} & Fixed IoT node, developing-area & No & No & Static \\
\citet{zhang2020realtime} & Mobile + fixed IoT network & Yes (short-term) & No & Static \\
\citet{kataria2022aiiot} & City-scale, network-assisted & Yes & No & Static \\
\citet{mani2021aipowered} & Fixed IoT node & Yes & No & Static \\
\citet{ramadhani2025iot} & Fixed IoT node, low-cost & No & No & Static \\
\citet{pavan2025realtime} & Fixed IoT node, alert-based & No & No & Static \\
\citet{shabbir2025paradigm} (review) & Low-cost sensor networks & Mixed & No & Static \\
\citet{dubey2021assessing} & Room-scale purifier study & No & Yes (reactive) & N/A \\
\citet{lu2023realworld} & Home-scale purifier study & No & Yes (reactive) & N/A \\
\citet{lee2022effect} & Vehicle-scale (bus) purifier & No & Yes (reactive) & N/A \\
\citet{chen2025fine} & Home-scale, wildfire event & No & Yes (reactive) & N/A \\
\textbf{NirmalNode (proposed)} & \textbf{Localized hotspot node} & \textbf{Yes (1--2 hr ahead)} & \textbf{Yes (predictive)} & \textbf{Incremental / Adaptive} \\
\bottomrule
\end{tabular}
\end{table}

\section{Research Gap}
\label{sec:lr-research-gap}

Table~\ref{tab:lr-gap} consolidates the specific gaps identified across the preceding sections into the concrete design requirements that motivate NirmalNode.

\begin{table}[htbp]
\centering
\caption{Identified research gaps and corresponding NirmalNode design response}
\label{tab:lr-gap}
\small
\begin{tabular}{@{}p{4.3cm}p{4.0cm}p{4.6cm}@{}}
\toprule
\textbf{Gap Identified in Literature} & \textbf{Representative Sources} & \textbf{NirmalNode Design Response} \\
\midrule
Monitoring/prediction not coupled to an automated physical mitigation action & \citep{alam2025iot,zhang2020realtime,kataria2022aiiot,mani2021aipowered} & Prediction output directly drives purifier activation (Decision Engine, Ch. 3) \\
Purification evaluated only at room/vehicle scale, not worker-proximate hotspot scale & \citep{dubey2021assessing,lu2023realworld,lee2022effect,chen2022efficacy,chen2025fine} & Compact cyclone--HEPA--activated-carbon unit sized for a single hotspot / worker zone \\
Purifier activation is reactive (responds to current conditions only) & \citep{dubey2021assessing,lu2023realworld,lee2022effect} & 1--2 hour ahead predictive activation before hazard threshold is crossed \\
AI models are trained once and deployed statically; no post-deployment adaptation & \citep{alam2025iot,zhang2020realtime,kataria2022aiiot,mani2021aipowered} & Incremental/online learning (SGD, PA, Hoeffding Tree, ARF) for continuous, area-specific model update \\
Health/exposure evidence establishes hotspot severity but proposes no engineering countermeasure & \citep{nasri2023pm25,sekhavati2023particulate,alhadhrami2024pm25mental,rahman2025perceived} & Hotspot-targeted purification directly engineered around documented worker exposure points \\
Low-cost sensor deployment is necessary but insufficient without an integrated mitigation layer & \citep{shabbir2025paradigm,karmakar2024exploring,noman2025aiair} & Full sense--predict--purify pipeline built entirely on low-cost, locally available components \\
Very low real-world adoption of purification technology in Bangladesh due to cost and misbeliefs & \citep{chowdhury2025misbeliefs} & Low unit cost, localized (not whole-room) purification volume, and demonstrable predictive operation intended to reduce both cost and trust barriers \\
\bottomrule
\end{tabular}
\end{table}

\needspace{8\baselineskip}
\section{Chapter Summary}
\label{sec:lr-summary}

In this chapter we reviewed four intersecting bodies of literature – ambient and industrial air pollution in Bangladesh, IoT-based air quality monitoring, AI-based pollution prediction and air purification technology – and have shown that while each individual component of NirmalNode (low-cost IoT sensing, AI-based prediction, HEPA/activated-carbon/cyclone purification) is independently well validated in prior work, no reviewed system integrates all of them into a single hotspot-targeted, predictively-activated, incrementally learning purification unit. The integration of this is the main novelty of the system presented in the rest of this capstone, in particular the integration of incremental/online learning for area-specific adaptation. NirmalNode describes the full methodology by which this combination is realized in Chapter 3.